\chapter*{قبل أن نبدأ}

هذا الكتاب عن آلة تزن كيلوغرامًا ونصفًا، وتستهلك عشرين واطًا كمصباح خافت، ومع ذلك تتعرف على الوجوه، وتتعلم اللغات، وتلتقط الكرة وهي تطير. إنه عن دماغك. سننظر إليه كما ينظر المهندس إلى لوحة إلكترونية لا يعرفها: ما القطع، وكيف وُصلت، وأي إشارات تجري في الأسلاك، وماذا تحسب هذه الدائرة.

لن تحتاج إلى معادلات ولا إلى معرفة بالأحياء. يكفي الفضول والاستعداد لتخيّل سلك ومصباح ومفتاح. وكل ما عدا ذلك سنجمعه في الطريق.

للكتاب نسخة ثانية كاملة، فيها المعادلات والنماذج والتمارين والإحالات إلى التجارب. فصول النسختين تسير بالترتيب نفسه وتحمل الأرقام نفسها. إذا شدّك فصل ما، فافتح نسخته الكاملة: ستجد هناك من أين جاءت كل عبارة وكيف تتحقق منها بنفسك. وعلى موقع الكتاب يمكن تبديل كل فصل بين ”مختصر“ و”كامل“ بزر واحد.

تنبيه واحد. معرفتنا بالدماغ أبعد ما تكون عن الاكتمال، وسأقول بصدق أين ينتهي المقيس ويبدأ التخمين. هذا ليس عيبًا في الكتاب، بل أمتع ما فيه: فصول كثيرة تنتهي بأسئلة لا يعرف أحد جوابها حتى الآن.
